% ===== PRÉAMBULE CANONIQUE — à coller VERBATIM en tête de CHAQUE .tex =====
\documentclass[11pt,a4paper]{article}
\usepackage[utf8]{inputenc}\usepackage[T1]{fontenc}\usepackage{lmodern}
\usepackage[french]{babel}
\usepackage{amsmath,amssymb,mathtools}\usepackage{icomma}
\usepackage[version=4]{mhchem}          % \ce{...} : équations & formules
\usepackage{siunitx}\sisetup{locale=FR} % \qty{}{}, \si{}, \ang{}
\usepackage{chemfig}                    % \chemfig{...} : structures développées
\usepackage{xcolor}\usepackage{colortbl}
\usepackage{tikz}\usepackage{pgfplots}\pgfplotsset{compat=1.18}
\usepackage[most]{tcolorbox}\usepackage{enumitem}\usepackage{array,booktabs}
\usepackage[a4paper,margin=2.2cm]{geometry}
\usetikzlibrary{arrows.meta,calc,angles,quotes,positioning,patterns,backgrounds,shapes.geometric}
\definecolor{primary}{RGB}{222,49,99}\definecolor{primarydark}{RGB}{150,22,55}
\definecolor{defcol}{RGB}{0,114,178}\definecolor{methcol}{RGB}{52,92,148}
\definecolor{excol}{RGB}{0,135,90}\definecolor{attncol}{RGB}{200,40,55}
\tcbset{boxbase/.style={enhanced,breakable,boxrule=0.9pt,arc=2.5pt,
  left=3mm,right=3mm,top=2.3mm,bottom=2.3mm,fonttitle=\bfseries,coltitle=white}}
% --- boîtes TUTORIEL (noms EXACTS, ne pas renommer) ---
\newtcolorbox{cle}[1][]{boxbase,colback=defcol!6,colframe=defcol,
  title={\textsc{À retenir}\ifx\relax#1\relax\relax\else\textmd{ : \itshape #1}\fi}}
\newtcolorbox{methode}[1][]{boxbase,colback=methcol!7,colframe=methcol,sharp corners,
  title={\textsc{Méthode}\ifx\relax#1\relax\relax\else\textmd{ --- \itshape #1}\fi}}
\newtcolorbox{exemplebox}[1][]{boxbase,colback=excol!5,colframe=excol,
  title={\textsc{Exemple}\ifx\relax#1\relax\relax\else\textmd{ : \itshape #1}\fi}}
\newtcolorbox{piege}[1][]{boxbase,colback=attncol!6,colframe=attncol,
  title={\textsc{Piège}\ifx\relax#1\relax\relax\else\textmd{ : \itshape #1}\fi}}
\newtcolorbox{remarque}[1][]{boxbase,colback=black!4,colframe=black!45,
  title={\textsc{Remarque}\ifx\relax#1\relax\relax\else\textmd{ : \itshape #1}\fi}}
% --- boîtes EXERCICE / CORRIGÉ (noms EXACTS, auto-comptées) ---
\newtcolorbox[auto counter]{exo}[1][]{boxbase,colback=primary!4,colframe=primary,
  title={\textsc{Exercice}~\thetcbcounter\ifx\relax#1\relax\relax\else\textmd{ --- \itshape #1}\fi}}
\newtcolorbox[auto counter]{corrige}[1][]{boxbase,colback=excol!4,colframe=excol,
  title={\textsc{Corrigé}~\thetcbcounter\ifx\relax#1\relax\relax\else\textmd{ --- \itshape #1}\fi}}
\newcommand{\R}{\mathbb{R}}\newcommand{\N}{\mathbb{N}}\newcommand{\Z}{\mathbb{Z}}\newcommand{\ds}{\displaystyle}
% (un \usepackage{fancyhdr}+en-tête est possible ; il est ignoré côté web)
% =========================================================================
\begin{document}

% THEME_TITLE: La mole, le nombre d'Avogadro et la masse molaire
\noindent
Une réaction chimique met en jeu un nombre colossal d'entités : on ne les compte jamais une à une,
on les regroupe en \textbf{paquets} de taille fixe, la \textbf{mole}. Tout ce thème consiste à
circuler entre trois grandeurs — un \emph{nombre d'entités} $N$, une \emph{quantité de matière} $n$
(en moles) et une \emph{masse} $m$ (en grammes) — à l'aide de deux « ponts » seulement.

\begin{cle}[la mole et le nombre d'Avogadro]
La \textbf{mole} (\si{\mole}) est l'unité de \textbf{quantité de matière}. Une mole contient
$N_{\mathrm{A}} = 6{,}02\times10^{23}$ entités élémentaires (le \textbf{nombre d'Avogadro}).
\[ \boxed{\,N = n\times N_{\mathrm{A}}\,}\qquad\Longleftrightarrow\qquad n = \dfrac{N}{N_{\mathrm{A}}}. \]
On \textbf{multiplie} par $N_{\mathrm{A}}$ pour aller des moles aux entités, on \textbf{divise} par $N_{\mathrm{A}}$ pour
faire l'inverse. Il faut \emph{toujours} préciser la nature de l'entité : une mole \emph{d'atomes},
\emph{de molécules} ou \emph{d'ions}.
\end{cle}

\begin{methode}[calculer une masse molaire $M$]
La \textbf{masse molaire} $M$ est la masse d'une mole ; elle s'exprime en \si{\gram\per\mole}.
Pour une molécule, on \textbf{additionne} les masses molaires atomiques de tous les atomes de la
formule (lues dans la classification périodique) :
\[ \boxed{\,M(\text{composé}) = \sum_i \nu_i\, M(\text{élément}_i)\,} \]
où $\nu_i$ est le nombre d'atomes de l'élément $i$ dans la formule brute.
\end{methode}

\begin{exemplebox}[masse molaire de l'acide sulfurique \ce{H2SO4}]
Avec $M(\ce{H})=1$, $M(\ce{S})=32$, $M(\ce{O})=16$ (en \si{\gram\per\mole}) :
\[ M(\ce{H2SO4}) = \underbrace{2\times 1}_{\ce{H}} + \underbrace{32}_{\ce{S}} + \underbrace{4\times 16}_{\ce{O}}
   = 2+32+64 = \SI{98}{\gram\per\mole}. \]
\end{exemplebox}

\begin{cle}[le pont masse $\leftrightarrow$ moles]
\[ \boxed{\,n = \dfrac{m}{M}\,}\qquad\Longleftrightarrow\qquad m = n\times M. \]
On \textbf{divise} une masse par $M$ pour obtenir des moles ; on \textbf{multiplie} des moles par
$M$ pour obtenir une masse. Vérification des unités : $\dfrac{\si{\gram}}{\si{\gram\per\mole}}=\si{\mole}$.
\end{cle}

\begin{center}
\begin{tikzpicture}[scale=1,
   box/.style={draw,rounded corners=3pt,align=center,font=\footnotesize,minimum height=1cm,minimum width=2.7cm,line width=0.9pt}]
  \node[box,draw=primary,fill=primary!8] (m) at (0,0) {Masse\\ $m$ (\si{\gram})};
  \node[box,draw=excol,fill=excol!8] (n) at (4.6,0) {Quantité\\ $n$ (\si{\mole})};
  \node[box,draw=defcol,fill=defcol!8] (N) at (9.2,0) {Nombre d'entités\\ $N$};
  \draw[-{Stealth[length=2mm]},line width=0.9pt] ([yshift=4pt]m.east) -- node[above,font=\scriptsize]{$\div M$} ([yshift=4pt]n.west);
  \draw[-{Stealth[length=2mm]},line width=0.9pt] ([yshift=-6pt]n.west) -- node[below,font=\scriptsize]{$\times M$} ([yshift=-6pt]m.east);
  \draw[-{Stealth[length=2mm]},line width=0.9pt] ([yshift=4pt]n.east) -- node[above,font=\scriptsize]{$\times N_{\mathrm{A}}$} ([yshift=4pt]N.west);
  \draw[-{Stealth[length=2mm]},line width=0.9pt] ([yshift=-6pt]N.west) -- node[below,font=\scriptsize]{$\div N_{\mathrm{A}}$} ([yshift=-6pt]n.east);
\end{tikzpicture}
\end{center}
\begin{center}\footnotesize\itshape La quantité de matière $n$ est le carrefour : tout passe par les moles.\end{center}

\begin{piege}[deux erreurs classiques]
\begin{itemize}[leftmargin=1.4em,topsep=1pt,itemsep=1pt]
  \item La masse molaire s'exprime en \si{\gram\per\mole}, \textbf{jamais en uma} (l'uma sert à peser
        \emph{un} atome, pas une mole).
  \item « Une mole \emph{de quoi} ? » : $\SI{1}{\mole}$ de \ce{CO2} contient $\SI{2}{\mole}$ d'atomes
        d'oxygène. On multiplie par l'indice de la formule.
\end{itemize}
\end{piege}

\end{document}
